%% metadata.tex
%% Central de Dados e Metadados do Trabalho Acadêmico
%% Altere as informações abaixo com os dados reais do seu projeto para preencher automaticamente
%% a Capa, Folha de Rosto, Folha de Aprovação e demais elementos pré-textuais.

% ==============================================================================
% 1. INFORMAÇÕES PESSOAIS E DO TRABALHO
% ==============================================================================
\newcommand{\meuautor}{Nome Completo do Autor}
\newcommand{\meutitulo}{Título Completo do Trabalho Acadêmico}
\newcommand{\meusubtitulo}{Subtítulo do Trabalho (se aplicável)} % Se não houver, deixe em branco
\newcommand{\meulocal}{Cidade -- UF}
\newcommand{\meuano}{2026}

% ==============================================================================
% 2. INFORMAÇÕES DA INSTITUIÇÃO
% ==============================================================================
\newcommand{\minhainstituicao}{Nome da Universidade ou Faculdade}
\newcommand{\meudepartamento}{Centro de Ciências Exatas e Tecnológicas}
\newcommand{\meucurso}{Curso de Graduação em Nome do Curso}

% ==============================================================================
% 3. INFORMAÇÕES ACADÊMICAS E DO TIPO DE TRABALHO
% ==============================================================================
% Ex: Trabalho de Conclusão de Curso (TCC), Dissertação, Tese, Relatório Técnico
\newcommand{\tipotrabalho}{Trabalho de Conclusão de Curso (TCC)} 

% Grau acadêmico pretendido (Ex: Bacharel em Ciência da Computação, Mestre em Física)
\newcommand{\grauacademicopretendido}{Bacharel em Ciência da Computação}

% Texto de apresentação para a Folha de Rosto e Aprovação
% Este texto varia conforme o tipo de trabalho e grau. Ajuste com os seus dados.
\newcommand{\textopresentacao}{%
  Trabalho de Conclusão de Curso apresentado ao \meucurso\ do \minhainstituicao, 
  como requisito parcial para obtenção do título de \grauacademicopretendido.
}

% ==============================================================================
% 4. ORIENTAÇÃO E COORDENAÇÃO
% ==============================================================================
\newcommand{\meuorientador}{Prof. Dr. Nome do Orientador}
\newcommand{\titulacaoorientador}{Orientador} % Ex: Orientador

\newcommand{\meucoorientador}{Prof. Msc. Nome do Coorientador} % Se não houver, comente nos templates
\newcommand{\titulacaocoorientador}{Coorientador}

% ==============================================================================
% 5. BANCA EXAMINADORA (Para a Folha de Aprovação)
% ==============================================================================
\newcommand{\dataaprovacao}{Data de Defesa}

\newcommand{\avaliadorum}{Prof. Dr. Nome do Primeiro Avaliador}
\newcommand{\instituicaoavaliadorum}{\minhainstituicao\ -- \minhainstituicao}

\newcommand{\avaliadordois}{Prof. Dr. Nome do Segundo Avaliador}
\newcommand{\instituicaoavaliadordois}{Outra Instituição de Ensino -- SIGLA}
